\chapter{Cmd, Sub, and the runtime}
\label{ch:cmdsub}

Chapter~\ref{ch:elm} introduced \code{Cmd<Msg>} as the type of
``effects as data.'' This chapter goes deeper: every alternative of
\code{Cmd}, every alternative of \code{Sub}, how the runtime
interprets them, and how background work plays with the pure
\code{update} function.

\section{The full \code{Cmd<Msg>} variant}

\begin{lstlisting}
template <typename Msg>
class Cmd {
public:
    struct None        {};
    struct Quit        {};
    struct Batch       { std::vector<Cmd> cmds; };
    struct After       { std::chrono::milliseconds delay; Msg msg; };
    struct SetTitle    { std::string title; };
    struct WriteClipboard { std::string content; };
    struct Task        { std::function<void(std::function<void(Msg)>)> run; };
    struct IsolatedTask { std::function<void(std::function<void(Msg)>)> run; };
    struct CommitRows  { int rows; };
    // ... and a few more
};
\end{lstlisting}

Each alternative is a small struct. Smart constructors give the
ergonomic API:

\begin{lstlisting}
Cmd<Msg> cmd1 = Cmd<Msg>::quit();
Cmd<Msg> cmd2 = Cmd<Msg>::after(std::chrono::seconds(1), Tick{});
Cmd<Msg> cmd3 = Cmd<Msg>::write_clipboard("hello");
Cmd<Msg> cmd4 = Cmd<Msg>::batch({cmd1, cmd2, cmd3});
\end{lstlisting}

Let us tour each one.

\subsection{\code{None} --- the identity}

The do-nothing effect. Returned by \code{update} when there is no
side effect to perform. Equivalent to \code{Cmd<Msg>\{\}} (default
construction).

\subsection{\code{Quit} --- exit the program}

The runtime, on seeing \code{Quit}, breaks out of its main loop,
restores cooked mode, exits alt-screen, and returns from
\code{run<P>()}. \code{main} then returns.

\subsection{\code{Batch} --- multiple effects}

\begin{lstlisting}
Cmd<Msg> save_and_quit = Cmd<Msg>::batch({
    Cmd<Msg>::write_clipboard(model.text),
    Cmd<Msg>::quit(),
});
\end{lstlisting}

Batched commands are interpreted in order. The runtime walks the
vector and dispatches each one. There is no concurrency guarantee
\emph{within} a batch --- \code{Task} commands inside a batch may
finish in any order.

\subsection{\code{After} --- one-shot timer}

\begin{lstlisting}
return {model, Cmd<Msg>::after(std::chrono::seconds(3), Hide{})};
\end{lstlisting}

The runtime schedules a one-shot timer; when it fires, it dispatches
the \code{Hide\{\}} \code{Msg} into \code{update} as if it had come
from the user. Useful for toasts (auto-hide after N seconds), debouncing,
retry backoff.

\subsection{\code{SetTitle} --- set the terminal title}

Writes \code{ESC]0;\emph{title}ESC\textbackslash} to set the
terminal's window title. Most emulators honour it.

\subsection{\code{WriteClipboard} --- OSC-52 clipboard}

Writes \code{ESC]52;c;\emph{base64}ESC\textbackslash}, the OSC-52
``set system clipboard'' sequence. Modern emulators (kitty, iTerm2,
WezTerm) honour it, even over SSH. \maya{} provides a portable
path that falls back to xclip / pbcopy / clip.exe on terminals
that do not.

\subsection{\code{Task} --- async work on a shared pool}

The escape hatch for \emph{arbitrary} side effects:

\begin{lstlisting}
return {model, Cmd<Msg>::task([](auto dispatch) {
    auto result = http_get("https://example.com/data");
    dispatch(Loaded{std::move(result)});
})};
\end{lstlisting}

The lambda receives a \code{dispatch} callback. It runs on \maya{}'s
shared, elastic-but-bounded background thread pool. When it calls
\code{dispatch(msg)}, the runtime feeds \code{msg} into the next
\code{update} call as if it had come from any other source.

Use \code{Task} for short tasks: HTTP calls, file reads, tiny
shell-outs. The pool is bounded; if you saturate it, subsequent
tasks queue.

\subsection{\code{IsolatedTask} --- async work on a dedicated thread}

Same shape as \code{Task}, but each \code{IsolatedTask} gets its own
detached \code{std::thread}. Use this for work that can wedge
indefinitely on a blocking syscall: slow filesystem mounts, network
freezes, sandboxed subprocesses that may hang. A wedged isolated
task leaks one thread; a wedged \code{Task} would consume a slot in
the shared pool permanently. The trade is some thread-construction
overhead per call ($\sim$100-300 microseconds).

\subsection{\code{CommitRows} --- inline-mode scrollback commit}

As described in Chapter~\ref{ch:inline}: tells the inline renderer
to mark the top \code{n} rows of the current frame as committed to
scrollback. The next frame starts below them.

\section{The \code{Cmd} functor}

If you build composite UIs, each sub-component has its own
\code{Msg} type. To embed a child component's \code{Cmd<ChildMsg>}
into a parent's \code{Cmd<ParentMsg>}, \code{Cmd} provides a
\code{map} (the categorical-functor \code{fmap}):

\begin{lstlisting}
Cmd<ChildMsg> child_cmd = child_update(...);
Cmd<ParentMsg> parent_cmd = child_cmd.map([](ChildMsg cm) -> ParentMsg {
    return ParentMsg::FromChild{std::move(cm)};
});
\end{lstlisting}

\code{map} walks every alternative; for those that carry a \code{Msg}
(\code{After}, \code{Task}, \code{IsolatedTask}), it wraps the
inner \code{Msg} with the provided function. The result is a new
\code{Cmd} the parent runtime can interpret.

This is exactly Elm's \code{Cmd.map}; it is what makes the
architecture compose to many nested components without leaking
their message types upward.

\section{The full \code{Sub<Msg>} variant}

Subscriptions describe event sources you want active right now:

\begin{lstlisting}
template <typename Msg>
class Sub {
public:
    struct Keys     { std::function<std::optional<Msg>(Key)>   f; };
    struct Mouse    { std::function<std::optional<Msg>(Mouse)> f; };
    struct Resize   { std::function<std::optional<Msg>(Size)>  f; };
    struct Paste    { std::function<std::optional<Msg>(std::string)> f; };
    struct Focus    { std::function<std::optional<Msg>(bool)>  f; };
    struct Tick     { std::chrono::milliseconds period; std::function<Msg()> f; };
    struct Batch    { std::vector<Sub> subs; };
    // ...
};
\end{lstlisting}

Each subscription is a function that takes a typed event and
returns a \code{Msg} (or \code{nullopt} to ignore). The runtime
samples them on every iteration.

Builders:

\begin{lstlisting}
auto subs = Sub<Msg>::batch({
    key_map<Msg>({{'q', Quit{}}, {'r', Refresh{}}}),
    Sub<Msg>::tick(std::chrono::seconds(1), []{ return Tick{}; }),
    Sub<Msg>::mouse([](Mouse m) -> std::optional<Msg> {
        if (m.action == MouseAction::Press)
            return Click{m.x, m.y};
        return std::nullopt;
    }),
});
\end{lstlisting}

\code{subscribe} returns this whole tree on every frame, so it can
\emph{depend on the model}: a help-modal-open model returns
different keybindings than a help-modal-closed one. The runtime
re-samples each frame, so changing the subscription set is just
returning a different \code{Sub} from \code{subscribe}.

\section{The runtime loop, expanded}

Now we can sketch the runtime in full:

\begin{lstlisting}
template <Program P>
void run(RunConfig cfg) {
    auto [model, cmd0] = P::init();
    Terminal term{cfg.mode};                  // RAII raw mode + alt-screen
    Writer   writer{term.fd()};
    Renderer renderer{cfg};
    EventSource events{term.fd()};
    BgPool      bg;

    auto interpret = [&](Cmd<typename P::Msg> cmd) {
        std::visit(overload{
            [&](Cmd<Msg>::None) {},
            [&](Cmd<Msg>::Quit) { quit_requested = true; },
            [&](Cmd<Msg>::Batch& b) {
                for (auto& c : b.cmds) interpret(std::move(c));
            },
            [&](Cmd<Msg>::After& a) {
                bg.timer(a.delay, [a]{ events.dispatch(a.msg); });
            },
            [&](Cmd<Msg>::SetTitle& s) {
                writer.write(set_title(s.title));
            },
            [&](Cmd<Msg>::Task& t) {
                bg.run([t](auto dispatch) { t.run(dispatch); });
            },
            // ... more arms ...
        }, cmd);
    };

    interpret(std::move(cmd0));

    while (!quit_requested) {
        renderer.paint(P::view(model), writer);

        auto subs = P::subscribe(model);
        Event ev = events.wait();             // poll(2) + wake_fd

        auto msg = match_event(ev, subs);     // produce a Msg via subs
        if (!msg) continue;                    // event ignored

        auto [next_model, next_cmd] = P::update(std::move(model), *msg);
        model = std::move(next_model);
        interpret(std::move(next_cmd));
    }
}
\end{lstlisting}

That is the entire shape. Roughly 70 lines of real code; the
\maya{} version is longer because it handles error recovery,
inline-mode redraw, signal-driven dirty checking, and Windows.

\section{Other entry points}

\code{run<P>()} is the canonical entry point. \maya{} provides three
others for common short-form cases:

\begin{itemize}[leftmargin=*]
  \item \code{run(event\_fn, render\_fn)} --- the quick-tool form
    shown in Chapter~\ref{ch:intro}. Internally constructs a
    trivial \code{Program} that closes over signals.
  \item \code{live(render\_fn)} --- animated output, no event loop.
    Useful for progress bars and live monitors. Calls
    \code{render\_fn(dt)} every frame.
  \item \code{canvas\_run(setup, event\_fn, paint\_fn)} --- imperative
    Canvas access. Useful for games, demos, and visualisations
    where the abstraction of layout is in the way.
  \item \code{print(element)} --- one-shot, no event loop. Useful
    for command-line tools that just want a styled banner.
\end{itemize}

All four lead to the same underlying terminal driver and rendering
pipeline. The difference is in how state is managed.

\section{Testing a Program without a terminal}

Because \code{Program} is just four pure functions, you can drive
it without ever opening a terminal:

\begin{lstlisting}
TEST_CASE("counter goes up on Inc and stays put on Quit") {
    auto [m0, _0] = Counter::init();
    REQUIRE(m0.count == 0);

    auto [m1, _1] = Counter::update(m0, Counter::Inc{});
    REQUIRE(m1.count == 1);

    auto [m2, _2] = Counter::update(m1, Counter::Quit{});
    REQUIRE(m2.count == 1);   // Quit shouldn't increment
}
\end{lstlisting}

You can also assert on the \code{Cmd} that came back:

\begin{lstlisting}
auto [_, cmd] = Counter::update(m, Counter::Quit{});
REQUIRE(std::holds_alternative<Cmd<Msg>::Quit>(cmd.inner));
\end{lstlisting}

This is the testing superpower of effects-as-data: you don't need a
terminal, you don't need a clipboard, you don't need a network ---
you just match on the value the function returned.

\begin{recap}
\code{Cmd<Msg>} is the algebra of effects the runtime knows how to
perform. \code{Sub<Msg>} is the algebra of event sources the runtime
samples. Both are sum types with smart constructors and a
functor-style \code{map}. The runtime is a tiny loop: paint, wait
for an event, dispatch through the subs to get a \code{Msg}, call
\code{update}, interpret the returned \code{Cmd}, repeat.
\end{recap}

\section*{Workshop}

\begin{enumerate}[leftmargin=*]
  \item Write a toast: a \code{Msg} that shows a message, scheduled
    by \code{Cmd::after} to dispatch a \code{HideToast} after 3
    seconds.
  \item Add a \code{Cmd::task} that performs a slow operation
    (\code{std::this\_thread::sleep\_for(2s)}) and dispatches a
    \code{Loaded} \code{Msg}. Notice the UI stays responsive.
  \item Read \file{maya/include/maya/core/cmd.hpp}. Find the
    \code{map} method on \code{Cmd}. Trace how it walks the
    variant and rewraps the inner \code{Msg}.
\end{enumerate}
